\section{From packet capacity to weighted tubes}\label{sec:density}

We first obtain a positive tube estimate from the Kakeya maximal theorem.
We then use a decomposition by capacity to apply Theorem~\ref{thm:packet}
to an arbitrary initial array without imposing a condition on the sizes
of its coefficients. Throughout this section, the frequency and position
label bounds are fixed as in the packet setup. The positive
tube-duality formulation follows the classical mixed-norm framework
of Tao, Vargas, and Vega \cite[arXiv version~1, Section~3, Equation~(21)]{TaoVargasVega1998};
the capacity decomposition below is proved here.

For an initial array \(d=(d_{\nu,z})\), define
\begin{equation}\label{eq:tube-function}
\begin{split}
 T_{\nu,z}
   &=\{(X,t)\in\mathbb R^2\times\mathbb R:
                   0\le t\le1,\ |X-z-tU_\nu|\le\delta\},\\
 H_d(X,t)&=\sum_{\nu,z}|d_{\nu,z}|^2\mathbf1_{T_{\nu,z}}(X,t),
 \qquad b_\nu(d)=\sum_z|d_{\nu,z}|^2.
\end{split}
\end{equation}
All integrals of tube functions use Lebesgue measure \(dX\,dt\).
In particular, \(|T_{\nu,z}|=\pi\delta^2\) and
\begin{equation}\label{eq:tube-mass}
 \int_{\mathbb R^3}H_d=\pi m(d).
\end{equation}

For clarity, the precise maximal operator used below is
\[
 K_\rho g(\omega)=\sup_{a\in\mathbb R^3}
       \frac{1}{\pi\rho^2}\int_{\mathcal T_\rho(a,\omega)}|g(y)|\,dy,
 \qquad
 \mathcal T_\rho(a,\omega)
 =\{a+r\omega+w: |r|\le\tfrac12,\ w\perp\omega,\ |w|\le\rho\}.
\]
Thus \(\mathcal T_\rho(a,\omega)\) is a circular cylinder of length one
and volume \(\pi\rho^2\). The companion theorem
\cite[Theorem~1.1]{OpenAIKakeya2026} states that for every \(s>0\) there is a finite
constant \(M(s)\), which we enlarge to be at least one, such that
for every \(g\in L^3(\mathbb R^3)\),
\begin{equation}\label{eq:kakeya-input}
 \|K_\rho g\|_{L^3(S^2)}
 \le M(s)\rho^{-s}\|g\|_{L^3(\mathbb R^3)},
 \qquad 0<\rho<1.
\end{equation}
Here \(S^2\) carries its usual surface measure. We will use only
\(0<s\le1\).

\Needspace{5\baselineskip}
\begin{lemma}[Weighted Kakeya estimate]\label{lem:weighted}
There are fixed constants \(C\) and \(N_0\) such that for every dyadic
\(N\ge N_0\), every initial array \(d\), and every \(0<s\le1\),
\begin{equation}\label{eq:weighted-tubes}
 \int H_d^{3/2}
 \le C M(s)^{3/2}\delta^{-3s/2}
                 \delta^2\sum_\nu b_\nu(d)^{3/2}.
\end{equation}
The constants \(C,N_0\) depend only on the fixed geometric setup.
\end{lemma}

\begin{proof}
By Equation~\eqref{eq:velocity-monotonicity}, the velocity map is
bi-Lipschitz on the fixed frequency-label range, where it is bounded. The map
\(U\mapsto (U,1)/\sqrt{1+|U|^2}\) is bi-Lipschitz on this bounded
range: its inverse is \(\omega\mapsto\omega'/\omega_3\), whose
derivative is bounded where \(\omega_3\) is bounded away from zero.
Thus the directions
\[
 \omega_\nu=\frac{(U_\nu,1)}{\sqrt{1+|U_\nu|^2}}
\]
are separated by a fixed positive multiple of \(\delta\).
Choose spherical caps \(\Omega_\nu\) centered at \(\omega_\nu\)
with radius \(c_1\delta\), where the fixed \(c_1>0\) is small enough
that these caps are disjoint. Their surface areas satisfy
\(c_2\delta^2\le |\Omega_\nu|\le C_2\delta^2\).

The center segment of \(T_{\nu,z}\) has Euclidean length
\(\sqrt{1+|U_\nu|^2}\), bounded by a fixed constant. For each
\(\omega\in\Omega_\nu\), its distance from the line through
\((z,0)\) in direction \(\omega\) is \(O(\delta)\). Adding the
horizontal cross-section of radius \(\delta\) preserves this bound.
The projection of the entire tube onto that line is an interval of
bounded length. Cover that interval by a bounded number of intervals of
length one. It follows that \(T_{\nu,z}\) is contained in at most
\(C_3\) cylinders \(\mathcal T_{C_4\delta}(a,\omega)\), where
\(C_3\) and \(C_4\ge1\) are fixed. In particular, for every
\(g\in L^3(\mathbb R^3)\),
\[
 \int_{T_{\nu,z}}|g|
 \le C\delta^2 K_{C_4\delta}g(\omega)
 \qquad (\omega\in\Omega_\nu).
\]
Averaging over \(\Omega_\nu\), multiplying by \(|d_{\nu,z}|^2\),
and summing gives
\[
 \int H_d|g|
 \le C\sum_\nu b_\nu(d)
                  \int_{\Omega_\nu}K_{C_4\delta}g(\omega)\,d\omega.
\]
Set \(B_d=\sum_\nu b_\nu(d)\mathbf1_{\Omega_\nu}\). Disjointness
of the caps gives
\[
 \|B_d\|_{L^{3/2}(S^2)}
 \le C\left(\delta^2\sum_\nu b_\nu(d)^{3/2}\right)^{2/3}.
\]
Choose \(N_0\) so that \(C_4\delta<1\). H\"older's inequality on
\(S^2\) and Equation~\eqref{eq:kakeya-input} now yield
\[
 \int H_d|g|
 \le C M(s)\delta^{-s}
       \left(\delta^2\sum_\nu b_\nu(d)^{3/2}\right)^{2/3}\|g\|_3.
\]
We used \((C_4\delta)^{-s}\le\delta^{-s}\), so the displayed
constant is independent of \(s\). Duality between \(L^{3/2}\) and
\(L^3\), followed by taking the power \(3/2\), proves
Equation~\eqref{eq:weighted-tubes}. The sum over positions causes no
loss because \(K_\rho\) already takes the supremum over cylinder
locations.
\end{proof}

The next observation compares a single ellipse test with the tube mass.
Recall that \(E\) is a centered ellipse with semiaxes
\(L\ge a\ge\delta\), and
\[
 W_\kappa(E)=L^{1+\kappa}a^{1-\kappa}
            =La(L/a)^\kappa\ge La.
\]
If a nonzero coordinate subarray \(q\) of \(d\) is supported on
\(U_\nu\in u+E\), \(z\in v+E\), then for \(0\le t\le1\)
the horizontal support of \(H_q(\cdot,t)\) lies in
\[
 v+tu+E+tE+B(0,\delta)\subset v+tu+3E.
\]
Indeed, \(B(0,\delta)\subset E\), and the convexity and central
symmetry of \(E\) imply \(E+tE=(1+t)E\). Therefore
\(|\operatorname{supp}H_q|\le9\pi La\). Equation~\eqref{eq:tube-mass}
and H\"older's inequality give
\begin{equation}\label{eq:bundle-lower}
 \int H_q^{3/2}
 \ge \frac{\big(\int H_q\big)^{3/2}}
              {|\operatorname{supp}H_q|^{1/2}}
 \ge \frac{\pi}{3}\,
              m(q)\left(\frac{m(q)}{W_\kappa(E)}\right)^{1/2}.
\end{equation}
The bound is uniform in the eccentricity \(L/a\), in the independent
translations \(u,v\), and in \(\kappa>0\).

\begin{lemma}[Decomposition by capacity]\label{lem:density}
Let \(d\) be a nonzero finite initial array, and let
\(n=\#\operatorname{supp}d\). For every \(0<\kappa<1/10\), there
is a decomposition into coordinate subarrays with pairwise disjoint
supports,
\[
 d=\sum_{i=1}^{J}d^{(i)},\qquad J\le C\log(2+n),
\]
such that
\begin{equation}\label{eq:density-cost}
 \sum_{i=1}^{J}m(d^{(i)})\mathcal K_\kappa(d^{(i)})^{1/2}
 \le C\int H_d^{3/2}.
\end{equation}
Both constants are absolute and independent of the nonzero coefficient
sizes and of \(\kappa\).
\end{lemma}

\begin{proof}
Write \(D=\max_{\nu,z}|d_{\nu,z}|^2>0\). A disk of radius
\(\delta\), independently translated to a prescribed velocity and
position, is an admissible test with weight \(\delta^2\). On the
other hand, every test has \(W_\kappa(E)\ge\delta^2\). Hence
\begin{equation}\label{eq:capacity-range}
 D\le\mathcal K_\kappa(d)\le nD.
\end{equation}
Capacity is monotone under restriction to a subset of coordinates.

Put \(A=D/n^2\). The cost of any remaining array with capacity at
most \(A\) is controlled by the contribution of a largest original
coefficient to the tube integral.
Indeed, any coordinate subarray \(r\) with
\(\mathcal K_\kappa(r)\le A\) satisfies
\[
 m(r)\mathcal K_\kappa(r)^{1/2}
 \le \delta^2 nD\sqrt{D/n^2}=\delta^2D^{3/2}.
\]
An original entry whose squared modulus is \(D\) contributes
\(D\mathbf1_{T_{\nu,z}}\) to \(H_d\), whence
\(\int H_d^{3/2}\ge\pi\delta^2D^{3/2}\).
Thus the cost of any such remainder is controlled by
\(\pi^{-1}\int H_d^{3/2}\), independently of which entries remain.

Start with remainder \(r_1=d\).
Whenever \(\mathcal K_\kappa(r_j)>A\), choose an ellipse test whose
density is at least half of that supremum. Let \(q_j\) contain all
nonzero coordinates of \(r_j\) selected by that test, write
\(E_j\) for its ellipse, and put
\[
 \rho_j=\frac{m(q_j)}{W_\kappa(E_j)}
       \ge\tfrac12\mathcal K_\kappa(r_j),
 \qquad r_{j+1}=r_j-q_j.
\]
Such a test exists without assuming that the supremum is attained.
It selects at least one nonzero entry, so the process terminates after
at most \(n\) removals. Denote the terminal remainder by \(r_*\);
then \(\mathcal K_\kappa(r_*)\le A\), including the possibility
\(r_*=0\).

We group the removed bundles by the capacity of the remainder from
which they were selected. This controls the capacity of each group
by its first remainder, while leaving only logarithmically many groups.
For each integer \(k\ge0\), set \(\lambda_k=2^kA\), and group
together all \(q_j\) for which
\[
 \lambda_k\le\mathcal K_\kappa(r_j)<2\lambda_k.
\]
Let \(d_k\) be the union array at this level, omitting empty levels.
The array \(d_k\) is a coordinate restriction of the remainder at
the first removal in that level. Therefore
\(\mathcal K_\kappa(d_k)<2\lambda_k\). Every bundle at this level
has \(\rho_j\ge\lambda_k/2\), so additivity of \(m\) gives
\[
\begin{split}
 m(d_k)\mathcal K_\kappa(d_k)^{1/2}
 &\le \sqrt{2\lambda_k}\sum_{j\text{ at level }k}m(q_j)\\
 &\le 2\sum_{j\text{ at level }k}m(q_j)\sqrt{\rho_j}.
\end{split}
\]
Use Equation~\eqref{eq:bundle-lower} for each selected test. Since the
\(q_j\) are disjoint coordinate subarrays, their tube functions sum
to a function bounded by \(H_d\). For nonnegative numbers,
\(\sum_j x_j^{3/2}\le(\sum_jx_j)^{3/2}\). Consequently
\begin{equation}\label{eq:level-cost}
 \sum_km(d_k)\mathcal K_\kappa(d_k)^{1/2}
 \le C\sum_j\int H_{q_j}^{3/2}
 \le C\int H_d^{3/2}.
\end{equation}

The stopping estimate applies to \(r_*\). Adding its cost to
Equation~\eqref{eq:level-cost} proves
Equation~\eqref{eq:density-cost}.

Finally, at every removal the capacity lies in \((A,nD]\), by
monotonicity and Equation~\eqref{eq:capacity-range}. The ratio of these
endpoints is \(n^3\), so there are at most
\(1+\lfloor3\log_2 n\rfloor\) occupied dyadic levels. Including
the nonzero remainder gives \(J\le C\log(2+n)\). If \(n=1\),
then \(\mathcal K_\kappa(d)=D=A\); the process makes no removal
and the remainder alone gives the conclusion. No lower bound on the
smallest coefficient has been used.
\end{proof}

\Needspace{5\baselineskip}
\begin{proposition}[Discrete critical estimate]\label{prop:discrete}
Let \(c\) be the terminal coefficient array obtained from \(d\) by
the synthesis and sampling map of Theorem~\ref{thm:packet}, with any
fixed permitted output mask. For every sufficiently large dyadic \(N\),
\(0<\kappa<1/10\), \(0<\eta\le1\), and \(0<s\le1\), put
\(\gamma_0=10\kappa+\eta+3s/2\). Then
\begin{equation}\label{eq:discrete-critical}
 \sum_{\nu',z',j}|c_{\nu',z',j}|^3
 \le C N^6 P(\kappa,\eta)M(s)^{3/2}
             N^{\gamma_0}\log^2(2+N)
             \delta^2\sum_\nu b_\nu(d)^{3/2}.
\end{equation}
The constant \(C\) and the lower threshold for \(N\) are independent
of \(\kappa,\eta,s\). Their dependence is only on the fixed phase,
windows, and label bounds; all additional parameter dependence is
contained in \(P(\kappa,\eta)\) and \(M(s)\).
\end{proposition}

\begin{proof}
There is nothing to prove if \(d=0\). Otherwise use the decomposition
in Lemma~\ref{lem:density}, and let \(c^{(i)}\) be the terminal array
of \(d^{(i)}\) with the same output mask. The sampling map is linear,
so \(c=\sum_i c^{(i)}\). The triangle inequality in \(\ell^3\),
H\"older's inequality for the finite sum over \(i\), and
Theorem~\ref{thm:packet} give
\[
\begin{split}
 N^{-6}\sum|c|^3
 &\le J^2\sum_i N^{-6}\sum|c^{(i)}|^3\\
 &\le J^2P(\kappa,\eta)N^{10\kappa+\eta}
        \sum_i m(d^{(i)})\mathcal K_\kappa(d^{(i)})^{1/2}\\
 &\le C\log^2(2+n)P(\kappa,\eta)N^{10\kappa+\eta}
        \int H_d^{3/2}.
\end{split}
\]
There are \(O(N^2)\) initial frequency labels and \(O(N^2)\) initial
position labels in the fixed bounded ranges, so \(n\le C N^4\).
It follows that
\begin{equation}\label{eq:packet-tube-comparison}
 N^{-6}\sum|c|^3
 \le C P(\kappa,\eta)N^{10\kappa+\eta}\log^2(2+N)
                        \int H_d^{3/2}.
\end{equation}
Insert Lemma~\ref{lem:weighted} and use
\(\delta^{-3s/2}=N^{3s/2}\) to obtain
Equation~\eqref{eq:discrete-critical}. Every coordinate restriction
used above obeys the same label bounds as \(d\), so the same packet
constant \(P(\kappa,\eta)\) applies to every piece.
\end{proof}
